\documentclass[11pt]{article}
\usepackage[T1]{fontenc}
\usepackage[utf8]{inputenc}
\usepackage{lmodern,amsmath,amssymb,amsthm,microtype}
\usepackage[margin=25mm]{geometry}
\usepackage{hyperref}
\hypersetup{hidelinks,pdftitle={Completed algebraic closure descent for singular psh metrics}}
\setlength{\emergencystretch}{3em}
\newtheorem{theorem}{Theorem}
\newtheorem{lemma}[theorem]{Lemma}
\newtheorem{proposition}[theorem]{Proposition}
\newtheorem{corollary}[theorem]{Corollary}
\theoremstyle{remark}
\newtheorem{remark}[theorem]{Remark}
\newcommand{\PSH}{\operatorname{PSH}}
\newcommand{\MA}{\operatorname{MA}}
\newcommand{\an}{\mathrm{an}}
\newcommand{\ddc}{dd^c}
\newcommand{\E}{\mathcal E}
\newcommand{\Gal}{\operatorname{Gal}}
\title{Completed algebraic closure descent for singular psh metrics}
\author{P16.4 research continuation 214\\Unreviewed partial result}
\date{7 October 2026}
\begin{document}
\maketitle

\begin{abstract}
Let $K$ be any complete non-Archimedean field of characteristic zero,
including fields with nondiscrete value group, and let
$C=\widehat{\overline K}$. For a smooth geometrically integral projective
$K$-variety and an actual ample rational line bundle, we give an
all-point descent argument for Galois-invariant singular psh metrics
and their extended relative energy. The argument uses continuous
field-extension results and compact orbit maxima; it does not require
continuity of arbitrary psh envelopes. We also descend the
continuous-test energy dual by Haar averaging. These statements do not
construct the general unbounded mixed Monge--Ampere operator and do
not solve the entire original problem P16.4. All new statements in this
paper await independent verification.
\end{abstract}

\section{Setting and the precise result}

Fix a complete non-Archimedean real-valued field $K$ of characteristic
zero. Its absolute value may be trivial, discrete, or nondiscrete.
Let $X/K$ be smooth, geometrically integral and projective of dimension
$n\geq 1$, and let $L\in\operatorname{Pic}(X)\otimes\mathbb Q$ be an
actual ample rational line bundle. Set
\[
 T=X^{\an},\qquad C=\widehat{\overline K},\qquad
 Y=X_C^{\an},\qquad \pi:Y\longrightarrow T,\qquad
 G=\Gal(\overline K/K),\qquad V=(L^n)>0.
\]
The unique extension of the absolute value makes every element of $G$
an isometry of $\overline K$, hence of $C$. The induced action on $Y$
is semilinear. We use the profinite topology on $G$.

Choose a continuous psh reference metric $\psi$ on $L$, and put
$\psi_C=\pi^*\psi$. Such a reference exists: a generated positive
multiple of $L$ supplies a Fubini--Study metric. We write a metric as
$\psi+u$, so that $u$ is a scalar on the whole carrier. For brevity
$u\in\PSH_K(L;\psi)$ means that $\psi+u$ is a psh metric.
Here psh has exactly the definition in \cite[Definition 7.2]{BE}:
it is an all-point decreasing-net limit of continuous Fubini--Study
metrics on the fixed rational bundle. We exclude the identically
$-\infty$ metric; all other poles, including supported analytic points,
are retained. Constant additions are literal real constants. No torus,
skeleton, countability or support restriction is imposed.

For continuous psh $v$ define the raw relative energy
\begin{equation}\label{eq:energy}
 E_K(v)=\frac1{n+1}\sum_{j=0}^{n}
 \int_T v\,(\ddc(\psi+v))^j\wedge(\ddc\psi)^{n-j}\wedge\delta_X.
\end{equation}
The continuous mixed measures here are the actual measures of
\cite[Theorem 8.4]{BE}. Their total mass is $V$. Define, for singular
$u\in\PSH_K(L;\psi)$,
\begin{equation}\label{eq:extended}
 \begin{aligned}
 E_K(u)&=\inf\{E_K(v):v\in C(T)\cap\PSH_K(L;\psi),\ v\geq u\},\\
 \E_K^1&=\{u\in\PSH_K(L;\psi):E_K(u)>-\infty\}.
 \end{aligned}
\end{equation}
An approximating FS metric provides a continuous majorant, so the
infimum is well-defined in $[-\infty,\infty)$. Use the same definitions
over $C$ with reference $\psi_C$. This definition does not assume an
unbounded Monge--Ampere operator.

For a positive Radon measure $\mu$ on $T$ of raw mass $V$, set
\begin{equation}\label{eq:dual}
 J_K(\mu)=\sup_{v\in C(T)\cap\PSH_K(L;\psi)}
 \left(E_K(v)-\int_T v\,d\mu\right).
\end{equation}
Constants cancel in this expression. We call this the
\emph{continuous-test dual}; no equivalence with an arbitrary-field
all-$\E^1$ integrability condition is built into the definition.

\begin{theorem}\label{thm:main}
In the stated setting the following hold.
\begin{enumerate}
\item The action $G\times Y\to Y$ is jointly continuous and
$\pi$ identifies $T$ with the compact Hausdorff quotient $Y/G$.
\item Pullback gives an order-preserving bijection
\[
 \PSH_K(L;\psi)\ \simeq\ \PSH_C(L_C;\psi_C)^G.
\]
It is a bijection of actual metrics at every point, including their
$-\infty$ values, and identifies the continuous subspaces.
\item For every $u\in\PSH_K(L;\psi)$,
\[
 E_C(\pi^*u)=E_K(u),
\]
including the value $-\infty$. Thus pullback identifies $\E_K^1$
with the invariant part of $\E_C^1$.
\item Every positive Radon measure $\mu$ on $T$ has a unique invariant
positive Radon lift $\widetilde\mu$ on $Y$. The lift preserves raw mass,
is continuous for weak convergence, and satisfies
\[
 J_C(\widetilde\mu)=J_K(\mu)\quad\text{if }\mu(T)=V.
\]
\item For continuous psh slots $v_1,\ldots,v_n$,
\[
 \pi_*\MA_{C,c}(\pi^*v_1,\ldots,\pi^*v_n)
 =\MA_{K,c}(v_1,\ldots,v_n).
\]
The upstairs measure is the unique invariant lift of the downstairs
measure. The mass on both sides is $V$, with no field-degree factor.
\end{enumerate}
\end{theorem}

The theorem concerns fixed ample rational bundles and smooth
geometrically integral projective varieties. It makes no assertion for
arbitrary transcendental classes or nonprojective analytic spaces.
The energy equality is an extended scalar identity; it must not be
read as an unbounded mixed-product theorem.

\section{Continuous facts used, and their boundary}

We use the following exact continuous results of \cite{BE}.
Corollaries 7.5--7.6 give finite maximum, decreasing closure and uniform
closure of continuous psh metrics. Theorem 7.31 recognizes a
\emph{continuous} metric as psh after any complete scalar extension,
provided $X$ is geometrically reduced and the bundle is semiample.
Our $X,L$ satisfy these assumptions over both $K$ and $C$.
The complete proof of that theorem uses the $o(m)$ norm estimate in
Theorem 6.6 and the FS operator; it does not assume envelope continuity.
We do not replace it by an assertion of exact equality of extended
supnorms.

Theorem 7.14 supplies uniform smooth semipositive regularization of
continuous psh metrics. Proposition 8.3 and Theorem 8.4 supply the
actual continuous mixed measures, integration by parts, uniform-to-weak
continuity and scalar-extension pushforward. Definition 9.13 and
Proposition 9.14 supply continuous relative energy, monotonicity,
constant shift, the cocycle law, $V$-Lipschitz continuity and scalar
extension invariance. The normalization is raw, as in
\eqref{eq:energy}. The ordinary analytic facts used below are
projective analytification being compact Hausdorff, the seminorm
description on affine charts, and complete fields being henselian.

For completeness, real convexity of the continuous psh cone needs no
extra envelope theorem. A rational convex combination of two FS
metrics is an FS metric on the same rational bundle: clear
denominators, tensor all pairs of the generated section systems, and
add their weights. Uniform approximation by FS metrics proves the
same assertion for rational combinations of continuous psh metrics.
Approximating a real coefficient uniformly by rationals then gives
real convexity by Corollary 7.6.

\begin{lemma}\label{lem:concave}
The continuous energy $v\mapsto E_K(v)$ is concave. The same holds
over $C$. It is invariant under the semilinear $G$-action upstairs.
\end{lemma}
\begin{proof}
We give the direct smooth argument to identify the scope of this
property. For two smooth semipositive metrics $\phi_0,\phi_1$ on $L$
put $h=\phi_1-\phi_0$ and $\phi_t=\phi_0+th$. Multilinearity of
curvatures and integration by parts give the energy polynomial and
its derivatives
\[
 \frac d{dt}E(\phi_t,\psi)=\int h(\ddc\phi_t)^n\wedge\delta_X,
 \qquad
 \frac {d^2}{dt^2}E(\phi_t,\psi)
 =n\int h\,\ddc h\wedge(\ddc\phi_t)^{n-1}\wedge\delta_X.
\]
These formulas remain valid when $\psi$ is first approximated by
smooth semipositive metrics. With $T_t=(\ddc\phi_t)^{n-1}\wedge
\delta_X$, the smooth Stokes rule used in
\cite[Proposition 8.3(iii)]{BE} gives
\[
 \int h\,\ddc h\wedge T_t
 =-\int dh\wedge d^ch\wedge T_t\leq0.
\]
Here $dh\wedge d^ch$ is a positive real $(1,1)$ form with the
standard $d^c$ convention; the closed semipositive curvature factors
preserve positivity. Compactness eliminates boundary terms. Hence the
energy polynomial is concave on $[0,1]$. Theorem 7.14 and the uniform
continuity of mixed measures and energy extend the inequality to
continuous psh metrics and a continuous psh reference. This is the
continuous concavity also stated in \cite[Section 1.6]{BGM}; we use no
singular concavity statement from that source.

Naturality under a valued-field automorphism identifies the curvature
forms and analytic fundamental classes after the induced semilinear
isomorphism. Integration on the isomorphic carrier is unchanged.
Consequently $E_C(g^*v,g^*\psi_C)=E_C(v,\psi_C)$ for smooth metrics,
and then for continuous psh metrics by uniform regularization. Since
$g^*\psi_C=\psi_C$, the claimed invariance follows. The same argument
gives the equivariance of the continuous mixed operator.
\end{proof}

\section{The completed algebraic closure really is an orbit quotient}

\begin{lemma}\label{lem:quotient}
The action $G\times Y\to Y$ is jointly continuous. The fibers of
$\pi$ are nonempty and are exactly the $G$-orbits.
\end{lemma}
\begin{proof}
First let $E/K$ be finite Galois and $x\in T$. The analytic fiber over
$x$ is described by the multiplicative norms on
$\mathcal H(x)\otimes_K E$, where $\mathcal H(x)$ is the completed
residue field. This nonzero finite etale algebra is a product of finite
separable fields over the complete henselian field $\mathcal H(x)$.
Each field factor has a unique extension of its absolute value, so
the fiber is a finite nonempty set. The group $\Gal(E/K)$ acts
transitively on the factors: an orbit which were proper would give a
nontrivial invariant idempotent, whereas the invariant algebra is
$\mathcal H(x)$. Uniqueness of the absolute values also makes the
action transitive on the norm points. This includes nontrivial
valuations on $\mathcal H(x)$ and supported points of $T$.

We next justify the completion step, instead of assuming that every
coefficient in $C$ is algebraic. There is a natural map
\begin{equation}\label{eq:inverse}
 Y\longrightarrow\varprojlim_{E/K\ \mathrm{finite\ Galois}}X_E^{\an}.
\end{equation}
For a compatible point family choose an affine $K$-chart containing
its common downstairs scheme support, with coordinate generators
$t_1,\ldots,t_r$. The family defines a multiplicative seminorm on
$A\otimes_K\overline K$. For an element of $A\otimes_K C$, express
it as a polynomial in the $t_i$ with finitely many coefficients in
$C$ and approximate those coefficients by elements of $\overline K$.
For a difference polynomial supported on finitely many monomials,
the ultrametric inequality bounds its seminorm by
\[
 \max_\alpha |c_\alpha-c'_\alpha|
                   \prod_{i=1}^r|t_i|^{\alpha_i}.
\]
Thus approximants are Cauchy in the seminorm. More explicitly,
$\big||a|-|b|\big|\leq|a-b|$ ensures convergence of their numerical
norms. The bound proves independence of the chosen approximants.
It also proves independence of polynomial representatives: a relation
in the finitely generated ideal of the chart can be approximated by
finite linear combinations of its defining $K$-relations; its
seminorm is zero. Addition, multiplication, the ultrametric inequality
and restriction to the given absolute value on $C$ pass to limits.
We obtain a unique point on the $C$-chart. These constructions agree
on overlaps, since localization at a function of nonzero seminorm is
unique.

This proves that \eqref{eq:inverse} is a continuous bijection.
Its source is compact and its target is Hausdorff, so it is a
homeomorphism. The finite fibers over $x$ form an inverse system
with surjective transition maps: at each later finite level the same
finite-etale, complete-field argument applied over the earlier level
extends each norm point. Their compact inverse limit is nonempty,
which proves surjectivity of $\pi$.

For joint continuity of the action, functions with coefficients in
$\overline K$ have a finite-index open stabilizer in $G$.
Their evaluations therefore vary jointly continuously in $(g,y)$.
For coefficients in $C$, approximate by algebraic coefficients.
On an affine coordinate neighborhood where the finitely many $|t_i|$
are bounded, the displayed estimate is uniform in both $y$ and $g$:
every $g$ acts isometrically on the coefficient field and fixes the
$K$-coordinates. Uniform limits of these evaluation functions are
continuous. Such neighborhoods give the pointwise seminorm topology,
so the full action is jointly continuous.

Finally, for $y,z\in Y$ with $\pi(y)=\pi(z)$, the sets
\[
 F_E=\{g\in G:(gy)_E=z_E\}
\]
are nonempty closed subsets of the compact group $G$. Any finite
intersection is nonempty by passing to a finite Galois compositum and
using finite-level transitivity. The finite intersection property
gives a $g$ matching every finite level. The uniqueness in
\eqref{eq:inverse} gives $gy=z$. Orbits are contained in fibers by
definition, proving equality.
\end{proof}

Since a quotient by a compact group on a compact Hausdorff space is
compact Hausdorff, the induced continuous bijection $Y/G\to T$ is a
homeomorphism. Invariant continuous scalars therefore descend
continuously. Invariant upper semicontinuous scalars with values in
$[-\infty,\infty)$ descend upper semicontinuously: each closed upper
level set is compact, and its image is closed. This is an all-point
statement, independent of the scheme support of a point.

\section{Compact orbit maxima replace finite coefficient descent}

\begin{lemma}\label{lem:orbit}
For a continuous psh scalar $h$ on $Y$, the orbit maximum
\[
 h^G(y)=\max_{g\in G}h(gy)
\]
is continuous, invariant and psh. It is the pullback of a continuous
psh scalar on $T$.
\end{lemma}
\begin{proof}
Joint continuity and compactness of $G\times Y$ show that the maximum
exists and is continuous. To check psh, for every $\epsilon>0$ choose
a finite set $g_1,\ldots,g_N$ such that
\[
 \left\|h^G-\max_{1\leq j\leq N} h\circ g_j\right\|_\infty
 <\epsilon.
\]
Indeed uniform continuity gives, near each $g$, a group neighborhood
on which $|h(gy)-h(g'y)|<\epsilon$ uniformly in $y$; a finite
subcover suffices. Pullback by each semilinear isometry sends an FS
system to an FS system, and hence preserves psh. Finite maximum is
psh by Corollary 7.5 of \cite{BE}, and uniform closure of the
continuous cone proves that $h^G$ is psh.

The quotient identifies $h^G$ with a continuous downstairs scalar
$a$. The metric $\psi+a$ is continuous, and its pullback is psh.
Theorem 7.31 of \cite{BE}, with the actual complete extension $C/K$,
shows that $\psi+a$ is psh. No exact descent of individual weighted
FS section lists was required. In particular no assertion that
nonzero coefficients have norm one is used.
\end{proof}

\begin{proposition}\label{prop:singular}
Pullback gives the psh bijection in Theorem \ref{thm:main}(2).
\end{proposition}
\begin{proof}
Pullback of FS metrics and their decreasing limits gives one
direction, and surjectivity of $\pi$ gives injectivity and order
reflection. For the converse let $u_C$ be an invariant upstairs psh
scalar and choose its defining decreasing net of continuous FS
scalars $h_i\downarrow u_C$. By Lemma \ref{lem:orbit}, every $h_i^G$
is the pullback of a continuous downstairs psh scalar $a_i$, and
$a_i$ decreases.

We prove the exact endpoint without interchanging an arbitrary
supremum and infimum. Fix $y\in Y$ and a finite threshold $b>u_C(y)$.
The compact sets
\[
 Q_i=\{g\in G:h_i(gy)\geq b\}
\]
are decreasing with empty intersection, since $u_C(gy)=u_C(y)<b$.
Compactness and directedness imply $Q_i=\varnothing$ on a tail.
Thus $h_i^G(y)<b$ on that tail. Since $h_i^G\geq u_C$, it follows
that $\inf_i h_i^G(y)=u_C(y)$. If $u_C(y)=-\infty$, apply the same
argument to each finite $b$.

The quotient descends $u_C$ to a scalar $u$ on $T$. The $a_i$
decrease to $u$ at every point, so Corollary 7.5 and Lemma 7.4 of
\cite{BE} show that $u$ is psh in the actual decreasing-FS-net
definition. It is not identically $-\infty$ because $u_C$ is not.
For a continuous endpoint, quotient descent is continuous and the
same Theorem 7.31 gives the continuous statement. Conversely a
continuous downstairs endpoint has continuous pullback.
\end{proof}

\section{Extended energy descends without an unbounded product}

\begin{lemma}\label{lem:energy-net}
For any decreasing net of continuous psh scalars $v_i\downarrow u$
on $T$,
\[
 E_K(v_i)\downarrow E_K(u).
\]
The assertion includes the value $-\infty$ and holds over $C$ too.
\end{lemma}
\begin{proof}
By monotonicity and \eqref{eq:extended}, $E_K(v_i)\geq E_K(u)$.
Let $w$ be any continuous psh majorant of $u$, and fix
$\epsilon>0$. The compact sets
$\{x:v_i(x)-w(x)\geq\epsilon\}$ are decreasing with empty
intersection. On a tail $v_i\leq w+\epsilon$ everywhere.
Monotonicity and the constant-shift rule for continuous energy give
\[
 \inf_i E_K(v_i)\leq E_K(w)+V\epsilon.
\]
First let $\epsilon\downarrow0$, then take the infimum over all $w$.
This proves equality, even when that last infimum is $-\infty$.
There is no integration of a measurable directed net in this proof.
\end{proof}

\begin{proposition}\label{prop:energy}
For every downstairs psh scalar $u$,
$E_C(\pi^*u)=E_K(u)$.
\end{proposition}
\begin{proof}
Take a defining continuous FS net $v_i\downarrow u$. Its pullback
decreases to $\pi^*u$. Continuous scalar-extension invariance gives
$E_C(\pi^*v_i)=E_K(v_i)$ for each $i$. Apply Lemma
\ref{lem:energy-net} on both carriers to these very nets.
\end{proof}

One may instead start with arbitrary upstairs continuous majorants
and orbit-maximize them; Proposition \ref{prop:singular} and Lemma
\ref{lem:energy-net} give the same conclusion. Exact field descent
of each original FS list is unnecessary. The energy class here is the
actual scalar class defined by \eqref{eq:extended}, not an abstract
completion whose elements may lack potentials.

For two continuous downstairs psh references $\psi,\psi'$, the cocycle
identity and taking the same majorant infimum give
\[
 E_{\psi'}(\phi)=E_\psi(\phi)-E_\psi(\psi')
\]
also for singular $\phi$. The difference of the references is
continuous and bounded on the compact carrier. Consequently membership
in the extended finite-energy class is reference-independent; all
descent formulas are covariant under an actual change of reference.

\section{Invariant Radon lifts and the continuous-test dual}

\begin{lemma}\label{lem:radon}
Pushforward identifies invariant positive Radon measures on $Y$ with
positive Radon measures on $T$. Its inverse preserves mass and is
continuous for the weak topologies.
\end{lemma}
\begin{proof}
Let $dg$ be probability Haar measure on $G$. For $f\in C(Y)$ put
\[
 Af(y)=\int_G f(gy)\,dg.
\]
Joint continuity and compactness show that $Af$ is continuous and
invariant. It descends to a continuous function $\overline{Af}$ on
$T$. The positive bounded functional
$f\mapsto\int_T\overline{Af}\,d\mu$ has norm $\mu(T)$.
Riesz representation gives a positive Radon measure
$\widetilde\mu$ on $Y$. Haar invariance gives group invariance, and
$A(f\circ\pi)=f\circ\pi$ gives
$\pi_*\widetilde\mu=\mu$ and $\widetilde\mu(Y)=\mu(T)$.

To check uniqueness, approximate $Af$ uniformly in $y$ by finite
weighted group averages. Such approximation follows by choosing a
finite clopen partition of $G$ on each piece of which $f(gy)$ varies
by at most $\epsilon$, uniformly in $y$, and taking one representative
from each piece with its Haar weight. If $\lambda$ is invariant,
each such finite average has the same integral as $f$. Thus
$\int f\,d\lambda=\int Af\,d\lambda$, and the pushforward determines
the latter. The same formula against continuous tests proves weak
continuity of the lift. Weak continuity of pushforward is immediate.
\end{proof}

For a finite Galois level $E/K$ with a fiber of $r$ norm points, the
invariant fiber probability has weight $1/r$ at each point.
The number $r$ need not be $[E:K]$. This finite-fiber statement must
not be confused with a blanket division by a field degree. The
scalar-extension pushforward law for intersection measures retains
the geometric mass $V$.

\begin{proposition}\label{prop:dual}
If $\mu(T)=V$, then $J_C(\widetilde\mu)=J_K(\mu)$.
\end{proposition}
\begin{proof}
Pullback of continuous downstairs tests and continuous energy
invariance give $J_C(\widetilde\mu)\geq J_K(\mu)$.
For an arbitrary continuous upstairs psh scalar $h$, its Haar average
$Ah$ is a uniform limit of finite convex averages of the scalars
$h\circ g$. Real convexity and uniform closure of the continuous psh
cone show that $Ah$ is continuous psh. It is invariant, so Theorem
7.31 and the quotient descend it to a continuous psh scalar $a$.

Continuous concavity and $G$-invariance of energy give, for each
finite convex average, energy at least $E_C(h)$. Uniform continuity
of energy therefore gives
$E_C(Ah)\geq E_C(h)$. Invariance of the lift and the same finite
average approximation give
$\int Ah\,d\widetilde\mu=\int h\,d\widetilde\mu$.
Since $Ah=\pi^*a$, we obtain
\[
 E_C(h)-\int h\,d\widetilde\mu
 \leq E_K(a)-\int a\,d\mu\leq J_K(\mu).
\]
Taking the supremum over all continuous upstairs psh $h$ proves the
other inequality, including infinite dual values.
\end{proof}

For any single downstairs psh scalar $u$, the ordinary Radon
pushforward law gives
\[
 \int_Y |\pi^*u|\,d\widetilde\mu=
 \int_T |u|\,d\mu,
\]
as extended nonnegative integrals; these functions are Borel. Thus
the integrability of all downstairs $\E_K^1$ tests is equivalent to
integrability of all \emph{invariant} upstairs $\E_C^1$ tests.
We do not infer integrability of every noninvariant singular upstairs
test from this identity or from Proposition \ref{prop:dual}.

\section{Continuous measures and what remains unproved}

For continuous psh slots on the actual fixed bundle, Proposition
8.3(iv) and the uniform extension in Theorem 8.4 of \cite{BE} give
the pushforward statement in Theorem \ref{thm:main}(5). Naturality
under $g$ and invariance of the slots make the upstairs mixed measure
invariant. Lemma \ref{lem:radon} identifies it with the unique
invariant lift. Proposition 8.5 gives raw mass $V$ on each carrier.
These facts complete the proof of Theorem \ref{thm:main}.

There is a useful precise reduction for the remaining work.
\emph{If} an actual equivariant mixed Radon operator is constructed
on all of $(\E_C^1)^n$, its pushforward on pulled-back slots defines
an operator on $(\E_K^1)^n$. Genuine decreasing-continuity and
weighted integral identities for these slots transfer whenever they
are proved upstairs: pushforward is weakly continuous, and each
downstairs Borel test is its actual pullback. This conditional
observation is not the construction of that upstairs operator.

Likewise, if an upstairs equation has a solution for the invariant
lift and every upstairs solution is unique modulo constants, then
the normalized solution is invariant. For example impose
$\sup_Y u=0$ relative to $\psi_C$; $g^*u$ has the same normalization
and the same target. Equivariance and uniqueness imply $g^*u=u$.
Proposition \ref{prop:singular} then descends that actual solution.
This argument requires actual existence, equivariance and
every-solution uniqueness upstairs; it provides none of them.

The following original gaps therefore remain.
\begin{enumerate}
\item A genuine full unbounded mixed operator on arbitrary-field
finite-energy potentials, including approximation independence,
weighted limits, locality and comparison, has not been constructed.
\item Continuity of arbitrary envelopes over a nondiscrete field,
the required normalized potential compactness, orthogonality and
both signs of the outer first variation remain to be established
where needed for the full variational theorem.
\item The equivalence between the original intrinsic all-$\E^1$
target condition and a usable dual bound, actual variational
attainment and every-solution uniqueness are not proved here.
\end{enumerate}
A value group which is dense remains embedded and nondiscrete in
every valued extension. Passing to $C$ does not make it discrete.
The accepted trivial-valuation P16.5 theorem is not a theorem over
a nontrivially valued $C$, and was not used as a grant in this paper.
These valuation restrictions cannot be removed by renaming the
base field. No original problem is declared solved by continuation 214.

\section{Source and verification record}

All asserted new descent steps are written above. The primary
continuous proof units were read in the literal version
\cite{BE}: Definition 7.2, Lemma 7.4 and Corollaries 7.5--7.6;
Theorem 7.14; Theorem 6.6 and the FS laws used by the complete proof
of Theorem 7.31; Proposition 8.3, Lemma 8.6 and the complete
extension proof of Theorem 8.4; Definition 9.13 and Proposition
9.14. The smooth Stokes and positivity foundations are standard
analytic differential calculus as referenced in Proposition 8.3;
we do not claim a fresh reading of every recursively cited
foundation. The continuous concavity statement in
\cite[Section 1.6]{BGM} is a scope cross-check; its cited
trivial-field treatment is not substituted for an arbitrary-field
singular theorem.

Local retrieval hashes and exact application scopes are recorded in
the accompanying quote-free source ledger. They are provenance,
not independent mathematical verdicts. Third-party source PDFs are
retained locally and are not part of the publication. No active
review output was read and no hot feedback was sent to the two
already running frozen reviews when writing this separate draft.
An independent whole-paper correctness review is required before
adding it to the accepted research-paper count.

\begin{thebibliography}{9}
\bibitem{BE} S. Boucksom and D. Eriksson,
\emph{Norms, determinant of cohomology and Fekete points},
arXiv:1805.01016v3 (8 January 2021).
\url{https://arxiv.org/pdf/1805.01016v3}.
\bibitem{BGM} S. Boucksom, W. Gubler and F. Martin,
\emph{Differentiability of relative volumes over an arbitrary
non-Archimedean field}, arXiv:2004.03847v1 (8 April 2020).
\url{https://arxiv.org/pdf/2004.03847v1}.
\end{thebibliography}
\end{document}
